\usepackage{fontspec}
\usepackage{xeCJK}
\setmainfont{Latin Modern Roman}
\setsansfont{Latin Modern Sans}
\setmonofont{Latin Modern Mono}[Scale=MatchLowercase]
\setCJKmainfont{FandolSong-Regular.otf}[BoldFont=FandolSong-Bold.otf,ItalicFont=FandolKai-Regular.otf]
\setCJKsansfont{FandolHei-Regular.otf}[BoldFont=FandolHei-Bold.otf]
\setCJKmonofont{FandolFang-Regular.otf}
\usepackage[a4paper,inner=23mm,outer=23mm,top=24mm,bottom=24mm,headheight=15pt]{geometry}
\usepackage{amsmath,amssymb,amsthm,bm,mathtools}
\usepackage{graphicx,booktabs,tabularx,array,longtable}
\usepackage{enumitem}
\usepackage{float}
\usepackage{xcolor}
\definecolor{CourseNavy}{HTML}{17324D}
\definecolor{CourseTeal}{HTML}{147D85}
\definecolor{CoursePale}{HTML}{F0F6F8}
\definecolor{CourseGray}{HTML}{536171}
\usepackage{titlesec}
\usepackage{setspace}
\setstretch{1.20}
\setlength{\parindent}{2em}
\setlength{\parskip}{0.25em}
\setlength{\emergencystretch}{2em}
\setlist{nosep,leftmargin=2em,topsep=0.4em}
\setlist[enumerate,1]{label=\arabic*.}
\usepackage[most]{tcolorbox}
\newtcolorbox{coursebox}[1]{enhanced,breakable,colback=CoursePale,colframe=CourseTeal,
  boxrule=0.45pt,arc=1mm,left=3mm,right=3mm,top=2mm,bottom=2mm,
  title={#1},fonttitle=\bfseries\sffamily,coltitle=white,
  before skip=1em,after skip=1em}
\titleformat{\chapter}[display]{\sffamily\bfseries\color{CourseNavy}\huge}
  {\large 第\thechapter 讲}{0.6em}{}
\titlespacing*{\chapter}{0pt}{-10pt}{20pt}
\titleformat{\section}{\Large\sffamily\bfseries\color{CourseNavy}}{\thesection}{0.6em}{}
\titleformat{\subsection}{\large\sffamily\bfseries\color{CourseTeal}}{\thesubsection}{0.6em}{}
\titleformat{\subsubsection}{\normalsize\sffamily\bfseries}{\thesubsubsection}{0.6em}{}
\setcounter{secnumdepth}{2}
\setcounter{tocdepth}{1}
\theoremstyle{definition}
\newtheorem{definition}{定义}[chapter]
\newtheorem{example}{例题}[chapter]
\newtheorem{exercise}{习题}[chapter]
\theoremstyle{plain}
\newtheorem{theorem}{定理}[chapter]
\newtheorem{proposition}[theorem]{命题}
\newtheorem{lemma}[theorem]{引理}
\newtheorem{corollary}[theorem]{推论}
\theoremstyle{remark}
\newtheorem{remark}{说明}[chapter]
\renewcommand{\proofname}{证明}
\newcommand{\R}{\mathbb R}
\newcommand{\E}{\mathbb E}
\DeclareMathOperator{\Var}{Var}
\DeclareMathOperator{\Cov}{Cov}
\DeclareMathOperator{\diag}{diag}
\DeclareMathOperator{\tr}{tr}
\DeclareMathOperator{\rank}{rank}
\DeclareMathOperator*{\argmin}{arg\,min}
\DeclareMathOperator*{\argmax}{arg\,max}
\newcommand{\norm}[1]{\left\lVert#1\right\rVert}
\newcommand{\inner}[2]{\left\langle#1,#2\right\rangle}
\usepackage[numbers,sort&compress]{natbib}
\usepackage{url}
\usepackage[unicode,colorlinks=true,linkcolor=CourseNavy,citecolor=CourseTeal,urlcolor=CourseTeal,
  bookmarksnumbered=true,pdfauthor={冯恺睿},pdftitle={凸优化与非线性优化：2027讲义}]{hyperref}
\usepackage{bookmark}
\usepackage{fancyhdr}
\pagestyle{fancy}
\fancyhf{}
\fancyhead[L]{\small\sffamily\color{CourseGray}凸优化与非线性优化}
\fancyhead[R]{\small\sffamily\color{CourseGray}2027 讲义}
\fancyfoot[C]{\small\thepage}
\renewcommand{\headrulewidth}{0.3pt}
\fancypagestyle{plain}{\fancyhf{}\fancyfoot[C]{\small\thepage}\renewcommand{\headrulewidth}{0pt}}
\renewcommand{\contentsname}{目录}
\renewcommand{\bibname}{参考文献}
\renewcommand{\figurename}{图}
\renewcommand{\tablename}{表}
\newcommand{\coursecover}[3]{%
  \hypersetup{pageanchor=false}
  \begin{titlepage}
  \setlength{\parindent}{0pt}
  \thispagestyle{empty}
  \vspace*{15mm}
  {\sffamily\color{CourseTeal}\Large 2027 \quad COURSE NOTES\par}
  \vspace{12mm}
  {\sffamily\bfseries\color{CourseNavy}\fontsize{31}{42}\selectfont 凸优化与\\非线性优化\par}
  \vspace{7mm}
  {\sffamily\color{CourseGray}\Large Convex and Nonlinear Optimization\par}
  \vspace{13mm}
  {\color{CourseTeal}\rule{\textwidth}{1.1pt}\par}
  \vspace{8mm}
  {\sffamily\bfseries\color{CourseNavy}\LARGE #1\par}
  \vspace{6mm}
  {\large #2\par}
  \vspace{12mm}
  {\normalsize #3\par}
  \vfill
  {\large 主讲：冯恺睿\par}
  \vspace{3mm}
  {\small 同济大学／上海创智学院课程材料\par}
  \vspace{5mm}
  {\small\color{CourseGray}从最优性条件、下降方法与随机噪声，到神经网络的更新几何与参数化。\par}
  \end{titlepage}
  \hypersetup{pageanchor=true}
}
